\documentclass[10pt,twocolumn,letterpaper]{article}
\usepackage{cvpr}
\usepackage{times}
\usepackage{graphicx}
\usepackage{amsmath,amssymb}
\usepackage[utf8]{inputenc}
\usepackage{url}
\usepackage{hyperref}
\usepackage{booktabs}
\usepackage{enumitem}
\setlength{\columnsep}{0.2in}
\setlength{\parindent}{1.5em}
\setlength{\parskip}{0pt}
\usepackage{indentfirst}

% Compact title spacing
\makeatletter
\def\@maketitle{%
  \newpage
  \null
  \vskip -0.4in
  \begin{center}%
    {\Large\bfseries \@title \par}%
    \vskip 0.3em%
    {\large
      \lineskip .4em%
      \begin{tabular}[t]{c}%
        \@author
      \end{tabular}\par}%
  \end{center}%
  \par
  \vskip 0.2em}
\makeatother

% Tighter section spacing
\makeatletter
\renewcommand\section{\@startsection{section}{1}{\z@}%
  {5pt plus 2pt minus 2pt}{3pt plus 1pt minus 0pt}%
  {\large\bfseries}}
\renewcommand\subsection{\@startsection{subsection}{2}{\z@}%
  {4pt plus 2pt minus 1pt}{2pt plus 1pt minus 0pt}%
  {\normalsize\bfseries}}
\makeatother

\begin{document}

\title{An Agentic RAG System for Medical Literature Question Answering}
\author{Ruoyu Zhang \quad Xu Chen\\
Auburn University, Auburn, AL, USA\\
{\tt\small \{ruz0002, xzc0066\}@auburn.edu}
}
\maketitle

\vspace{-2pt}
\section{Abstract}

We present a clinical Retrieval-Augmented Generation (RAG) agent that answers medical questions grounded in authoritative textbook knowledge. Built on a LangGraph-based agentic architecture, the system ingests 340K+ text chunks from 18 medical textbooks and employs hybrid retrieval (dense semantic search plus BM25 sparse matching) with cross-encoder reranking, hierarchical parent-child chunking, and token-aware context compression to produce source-attributed answers. Unlike fixed RAG pipelines, our agent dynamically plans retrieval strategies, decomposes complex queries into sub-questions, and compresses accumulated context to stay within token budgets. Evaluated on 100 USMLE-style questions from the MedQA benchmark, the system achieves 80\% overall accuracy and 86\% source attribution rate after three optimization iterations, improving from an initial 65\% accuracy. We further analyze the system's behavior along three safety-critical dimensions: in-knowledge-base accuracy, source grounding fidelity, and out-of-knowledge-base boundary awareness.

\section{Introduction}

Medical question answering requires both broad domain knowledge and precise clinical reasoning. Large language models (LLMs) have demonstrated remarkable general-purpose capabilities, yet they remain prone to hallucination when applied to high-stakes medical domains where factual accuracy is critical~\cite{ji2023survey}. Retrieval-Augmented Generation (RAG)~\cite{lewis2020retrieval} addresses this limitation by grounding LLM outputs in retrieved evidence, but standard single-pass RAG pipelines lack the ability to iteratively refine their retrieval strategy based on intermediate results.

This project builds an \textbf{agentic RAG system} for medical literature question answering. The key insight is that medical questions often require multi-step reasoning: identifying the relevant clinical domain, retrieving evidence from multiple sources, and synthesizing a grounded answer with proper source attribution. Our system implements this through a LangGraph~\cite{langgraph} workflow where an LLM-based orchestrator dynamically decides when to search, what to search for, and when sufficient evidence has been gathered to produce a final answer.

Our system differs from the original project proposal in two important ways. First, we shifted from a multimodal Medical Visual Question Answering (Med-VQA) setting to a \textbf{text-only} medical QA setting. The original plan used the VQA-RAD dataset~\cite{lau2018dataset} with radiology images, but we found that image-based medical datasets are poorly suited for RAG evaluation because many questions can be answered from the image alone without external knowledge retrieval. By adopting the MedQA benchmark~\cite{jin2021disease} (USMLE-style questions), we obtain a more rigorous evaluation of retrieval quality and knowledge grounding. Second, the agent architecture has become significantly more complex than originally planned, incorporating query decomposition, hierarchical chunking, hybrid retrieval with reranking, and token-aware context compression---features that emerged as necessary during iterative development.

\section{Background and Related Work}

\textbf{Retrieval-Augmented Generation.} RAG~\cite{lewis2020retrieval} augments language model generation with retrieved passages from an external knowledge base, enabling the model to produce answers grounded in factual evidence rather than relying solely on parametric knowledge. The core idea---retrieve relevant documents, then condition generation on them---has become a standard paradigm for knowledge-intensive NLP tasks.

\textbf{Medical RAG Systems.} MedRAG~\cite{xiong2024benchmarking} benchmarks RAG for medicine by providing standardized corpora from medical textbooks and clinical resources. Their work establishes that retrieval-augmented approaches significantly outperform closed-book LLMs on medical benchmarks. We adopt their textbook corpus as our knowledge base, comprising 18 authoritative medical textbooks.

\textbf{Agentic Reasoning.} The ReAct framework~\cite{yao2023react} interleaves reasoning traces with actions, enabling LLMs to dynamically plan and execute multi-step tasks. This paradigm is particularly relevant for medical QA, where the agent must decide \emph{what} to retrieve and \emph{when} retrieval is sufficient before generating an answer. Our orchestrator node implements a variant of ReAct within the LangGraph framework.

\textbf{Dense-Sparse Hybrid Retrieval.} Combining dense embeddings~\cite{reimers2019sentence} with sparse lexical matching (BM25) has been shown to improve retrieval robustness, particularly for domain-specific queries where exact terminology matters~\cite{lin2021few}. Cross-encoder rerankers further refine retrieval precision by scoring query-document pairs jointly~\cite{nogueira2019passage}.

\textbf{MedQA Benchmark.} MedQA~\cite{jin2021disease} is a large-scale, open-domain medical question-answering dataset based on USMLE board exam questions. Each question is multiple-choice with clinical vignettes, requiring both factual recall and clinical reasoning.

\section{Approach}

Our system is a two-level LangGraph agent with a top-level orchestration graph and per-question agent subgraphs.

\subsection{Knowledge Base Construction}

We build the knowledge base from 18 medical textbooks provided by MedRAG~\cite{xiong2024benchmarking}, totaling 340K+ indexed text chunks. Table~\ref{tab:textbooks} shows the distribution.

\begin{table}[h]
\centering
\small
\setlength{\tabcolsep}{3pt}
\begin{tabular}{lrlr}
\toprule
\textbf{Textbook} & \textbf{Chunks} & \textbf{Textbook} & \textbf{Chunks} \\
\midrule
Harrison's (Int.\ Med.) & 32,628 & Janeway (Immuno.) & 4,852 \\
Schwartz (Surgery) & 14,349 & Ross (Histology) & 4,411 \\
Adams (Neurology) & 12,370 & Levy (Physiology) & 4,370 \\
Williams (Obstetrics) & 9,166 & Nelson (Pediatrics) & 4,260 \\
Novak (Gynecology) & 7,947 & DSM-5 (Psychiatry) & 4,057 \\
Katzung (Pharm.) & 7,356 & Gray's (Anatomy) & 3,017 \\
Alberts (Cell Bio.) & 7,070 & Lippincott (Biochem.) & 1,973 \\
Robbins (Pathology) & 5,297 & First Aid Step 2 & 1,369 \\
\bottomrule
\end{tabular}
\vspace{2pt}
\caption{Knowledge base: 18 medical textbooks with chunk counts. Two smaller books (First Aid Step 1: 850, Pathoma: 505) are omitted for space.}
\label{tab:textbooks}
\end{table}

\textbf{Hierarchical Parent-Child Chunking.} Documents are first split by markdown headers (H1--H3) into semantically coherent sections. Small sections (${<}2000$ tokens) are merged with neighbors; large sections (${>}4000$ tokens) are recursively split. Each resulting \textit{parent chunk} (2000--4000 tokens) is then divided into \textit{child chunks} of 500 tokens with 100-token overlap. Child chunks serve as retrieval units (high precision), while parent chunks provide full context for answer generation (high recall).

\subsection{Hybrid Retrieval with Reranking}

Each child chunk is indexed with both dense and sparse representations:
\begin{itemize}[nosep,leftmargin=*]
  \item \textbf{Dense:} \texttt{all-mpnet-base-v2}~\cite{reimers2019sentence} (768-dim sentence embeddings) for semantic similarity search.
  \item \textbf{Sparse:} BM25 via \texttt{Qdrant/bm25} for exact lexical matching, critical for medical terminology (drug names, abbreviations).
\end{itemize}

Both indexes are hosted in Qdrant, which combines scores for hybrid retrieval. The system over-fetches $3\times$ candidates, then applies a cross-encoder reranker (\texttt{ms-marco-MiniLM-L-6-v2}~\cite{nogueira2019passage}) to produce the final top-$k$ results. After ranking, parent chunks are retrieved via stored parent IDs to provide broader context.

\subsection{Agent Architecture}

\textbf{Top-Level Graph.} Given a user query, the system: (1)~summarizes conversation history (keeping the last 6 messages), (2)~rewrites and decomposes the query into up to 3 sub-questions using structured LLM output, (3)~routes unclear queries to a clarification node (with human-in-the-loop interrupt), and (4)~fans out parallel agent subgraphs for each sub-question via LangGraph's \texttt{Send()} API, then aggregates sub-answers.

\textbf{Agent Subgraph.} Each sub-question is handled by an independent agent loop: an \textit{orchestrator} node (Gemini 2.5 Flash with tool bindings) decides whether to call \texttt{search\_child\_chunks} or \texttt{retrieve\_parent\_chunks}, observes the results, and either continues retrieval or produces a final answer. The loop is bounded by hard limits: 8 tool calls and 10 iterations per sub-question.

\textbf{Context Compression.} As the agent accumulates retrieved documents, a token-aware compression mechanism monitors the context size. When the total exceeds a dynamic threshold $T$:
\begin{equation}
T = T_{\text{base}} + T_{\text{current}} \times \alpha
\end{equation}
where $T_{\text{base}} = 2000$ tokens, $T_{\text{current}}$ is the current compressed summary length, and $\alpha = 0.9$ is a growth factor, the system compresses all retrieved content into a structured summary preserving clinical facts, dosages, evidence grades, and source attributions, while tracking already-executed searches and retrieved parent IDs to prevent duplicate retrieval.

\textbf{Fallback Mechanism.} If the agent exhausts its iteration or tool call budget without producing a final answer, a fallback node synthesizes the best possible answer from accumulated context without further tool access.

\subsection{Source Attribution and Safety}

Every answer includes explicit source citations (textbook name and section). The system is instructed to refuse answering when retrieved evidence is insufficient, rather than hallucinate---a critical safety property for medical applications. The orchestrator prompt enforces strict grounding: all factual claims must originate from retrieved documents, though the agent is allowed to \textit{reason} over retrieved evidence (e.g., applying clinical reasoning to connect retrieved facts).

\section{Intermediate Results}

\subsection{Evaluation Setup}

We evaluate on 100 questions sampled from MedQA~\cite{jin2021disease}: 80 \textit{in-knowledge-base} (in-KB) questions whose topics are covered by our 18 textbooks, and 20 \textit{out-of-knowledge-base} (out-KB) questions on topics outside our corpus. This split enables evaluation along three dimensions:
\begin{itemize}[nosep,leftmargin=*]
  \item \textbf{In-KB Accuracy:} Correctness when relevant knowledge exists.
  \item \textbf{Source Grounding:} Whether answers include proper citations.
  \item \textbf{Boundary Awareness:} Whether the system appropriately refuses out-KB questions rather than hallucinating.
\end{itemize}

The LLM backbone is Google Gemini 2.5 Flash with temperature 0 for deterministic outputs.

\subsection{Results}

Table~\ref{tab:results} shows the progression across three optimization iterations.

\begin{table}[h]
\centering
\small
\begin{tabular}{lccc}
\toprule
\textbf{Metric} & \textbf{v1} & \textbf{v2} & \textbf{v3} \\
\midrule
Overall Accuracy & 65.0\% & 75.0\% & \textbf{80.0\%} \\
In-KB Accuracy & 66.7\% & 80.0\% & 77.5\% \\
Source Attribution Rate & 66.7\% & 73.3\% & \textbf{86.2\%} \\
Out-KB Refusal Rate & 40.0\% & 40.0\% & 15.0\% \\
Avg.\ Time / Question & 32.1s & 35.3s & \textbf{26.0s} \\
\bottomrule
\end{tabular}
\vspace{2pt}
\caption{Evaluation metrics across three optimization iterations on 100 MedQA questions (80 in-KB, 20 out-KB).}
\label{tab:results}
\end{table}

\textbf{v1 $\to$ v2:} We added cross-encoder reranking and lowered the retrieval score threshold from 0.7 to 0.4, which significantly improved in-KB accuracy (+13.3 pp) by surfacing more relevant documents. We also strengthened grounding instructions in the orchestrator prompt.

\textbf{v2 $\to$ v3:} We balanced the grounding rules to allow the agent to \textit{reason} over retrieved evidence (not just quote it), improved answer extraction for multiple-choice formatting, and optimized the retrieval pipeline, reducing average latency from 35.3s to 26.0s while further improving overall accuracy to 80\%.

\subsection{Analysis}

The 80\% overall accuracy demonstrates that our agentic RAG approach can effectively answer USMLE-style medical questions. The 86.2\% source attribution rate indicates strong grounding behavior. The relatively low out-KB refusal rate (15\%) reveals a key challenge: the agent tends to attempt answers even when its knowledge base lacks relevant content, producing plausible but potentially incorrect responses---an important failure mode for safety-critical medical applications.

Qualitatively, we observe that the agent successfully handles multi-step clinical reasoning. For example, on a question about cholesterol embolization following cardiac catheterization, the agent correctly identified the diagnosis by retrieving relevant passages from Harrison's Internal Medicine and Robbins Pathology, integrating findings about eosinophilia, livedo reticularis, and characteristic biopsy findings.

Common failure modes include: (1)~ethical and procedural questions where textbook knowledge is insufficient, (2)~questions requiring knowledge not covered by any of the 18 textbooks, and (3)~over-confident answers on out-KB questions where the agent should refuse.

{\small
\bibliographystyle{ieee_fullname}
\begin{thebibliography}{10}
\setlength{\itemsep}{0pt}
\setlength{\parsep}{0pt}
\bibitem{lewis2020retrieval} P.~Lewis et al., ``Retrieval-Augmented Generation for Knowledge-Intensive NLP Tasks,'' \textit{NeurIPS}, 2020.
\bibitem{xiong2024benchmarking} G.~Xiong et al., ``Benchmarking Retrieval-Augmented Generation for Medicine,'' \textit{ACL Findings}, 2024.
\bibitem{yao2023react} S.~Yao et al., ``ReAct: Synergizing Reasoning and Acting in Language Models,'' \textit{ICLR}, 2023.
\bibitem{jin2021disease} D.~Jin et al., ``What Disease does this Patient Have? A Large-scale Open Domain Question Answering Dataset from Medical Exams,'' \textit{Applied Sciences}, 2021.
\bibitem{lau2018dataset} J.~J.~Lau et al., ``A Dataset of Clinically Generated Visual Questions and Answers about Radiology Images,'' \textit{Scientific Data}, 2018.
\bibitem{ji2023survey} Z.~Ji et al., ``Survey of Hallucination in Natural Language Generation,'' \textit{ACM Computing Surveys}, 2023.
\bibitem{reimers2019sentence} N.~Reimers and I.~Gurevych, ``Sentence-BERT: Sentence Embeddings using Siamese BERT-Networks,'' \textit{EMNLP}, 2019.
\bibitem{nogueira2019passage} R.~Nogueira and K.~Cho, ``Passage Re-ranking with BERT,'' \textit{arXiv:1901.04085}, 2019.
\bibitem{langgraph} LangChain, ``LangGraph: Build Stateful Multi-Actor Applications with LLMs.'' \url{https://github.com/langchain-ai/langgraph}, 2024.
\bibitem{lin2021few} J.~Lin et al., ``Few-Shot Text Ranking with Meta Adapted Synthetic Weak Supervision,'' \textit{ACL}, 2021.
\end{thebibliography}
}

\end{document}
